\section{Lean verification and a path to complete internalization}
\label{sec:lean}

\subsection{What has actually been checked}

The archive contains an arithmetic core whose source was successfully
compiled. The file \texttt{lean/DiophantineCore.lean} was checked
with Lean 4.32.0, commit
\nolinkurl{8c9756b28d64dab099da31a4c09229a9e6a2ef35}. Its only import is
\texttt{Std}; it has no Mathlib dependency, no added axioms, no unfinished
proofs, and no use of \texttt{native\_decide}.  The complete compiler output
and a source hash are in \texttt{evidence/lean-build.txt} and
\texttt{evidence/lean-build.json}.

Thirteen named theorems have been checked:
\begin{enumerate}
\item \texttt{anyLE\_correct}: a Boolean scan of $0,\ldots,n$ succeeds
  exactly when a successful index exists in that interval.
\item \texttt{findLE\_sound} and \texttt{findLE\_none\_iff}: the bounded
  witness search returns a valid index, and reports failure exactly when
  the Boolean scan reports failure.
\item \texttt{noBooleanDescent}: an explicitly iterated Boolean condition
  cannot persist under a strict decrease of a natural-number rank.
\item \texttt{quadratic\_expansion}, \texttt{second\_identity}, and
  \texttt{second\_on\_conic}: the quadratic expansion and second-intersection
  identities for an arbitrary integral ternary quadratic form.
\item \texttt{circle\_identity} and \texttt{circle\_inverse\_identity}:
  the polynomial identities underlying the circle parametrization.
\item \texttt{elliptic\_triangle\_identity},
  \texttt{elliptic\_triangle\_area},
  \texttt{triangle\_curve\_certificate}, and
  \texttt{triangle\_on\_curve}: the cleared-denominator identities relating
  a rational right triangle of area $n$ to $y^2=x^3-n^2x$.
\end{enumerate}

The output of \texttt{\#print axioms} lists only
\texttt{propext} and \texttt{Quot.sound}; for
\texttt{elliptic\_triangle\_area} it lists only \texttt{propext}.
In particular, \texttt{Classical.choice} and \texttt{sorryAx} do not occur
in any of the thirteen reports.  The polynomial proofs use a short local
tactic which expands integer powers and products by ring laws, puts
monomials into a fixed order, and asks \texttt{omega} to check the resulting
linear integer identity.  This avoids importing a classical polynomial
normalization theorem merely for the sake of convenient automation.

These are \emph{ordinary Lean theorems about integers and natural numbers}.
They are not yet proof objects of the deeply embedded theory \iPRA.
For example, the source proof of \texttt{anyLE\_correct} uses an ambient
induction whose displayed proposition contains an existential quantifier.
Its object-theory version must instead use the quantifier-free scan
invariants described below.  The absence of classical choice in the
ambient proofs is useful evidence, but does not by itself establish the
induction restriction or produce an \iPRA\ derivation.

\subsection{The existing ProveIt interface}

The inspected ProveIt snapshot is commit
\texttt{b5a3ad037723c44d69a3d5b28a950c554327a989}.
The relevant declarations are in the namespace
\texttt{LeanProofs.ConciseIPRA}.  The file \texttt{PR.lean} defines:
\begin{itemize}
\item \texttt{PR k}, the syntax of $k$-ary primitive recursive function
  symbols, with zero, successor, projections, composition, and primitive
  recursion;
\item \texttt{PTm n}, terms in $n$ free variables;
\item \texttt{Fm PTm n}, first-order intuitionistic formulas;
\item \texttt{PRAAx}, the defining equations, successor axioms, and
  induction for formulas certified quantifier-free;
\item \texttt{PRA.qfInd}, the derived rule that constructs precisely such
  an induction proof.
\end{itemize}
The theorem \texttt{pra\_iff\_native} in
\texttt{NativeBiInterpretation.lean} has the exact form
\[
 (\Gamma\vdash[\mathtt{PRAAx}]\varphi)
 \quad\longleftrightarrow\quad
 \mathtt{NPrf}\bigl(\Gamma.\mathtt{map}\ \mathtt{NFm.ofPRA}\bigr)
                     (\mathtt{NFm.ofPRA}\ \varphi).
\]
Thus the principal completion target for each theorem of this article is
an actual term of
\[
 []\vdash[\mathtt{PRAAx}]\varphi_G,
 \qquad \varphi_G:\mathtt{Fm}\ \mathtt{PTm}\ 0,
\]
where $\varphi_G$ is its explicitly encoded arithmetic statement.
Once that term is constructed, applying the forward direction of
\texttt{pra\_iff\_native} produces the corresponding native-scalar proof.
The bridge is not a rule turning any true theorem of Lean into a theorem
of \iPRA; it takes an \iPRA\ derivation as its input.

\subsection{A concrete module plan}

The following dependency order separates implementation work from
mathematical gaps.  The human proofs of the requisite lemmas are given in
the preceding chapters; the tasks here are their syntactic realization.

\paragraph{1. Encodings and primitive recursive symbols.}
Create \path{DG/Encoding.lean} and \path{DG/PRSymbols.lean}.
Use the signed-integer, rational, vector, polynomial, and finite-list
encodings fixed in the foundations chapter.  Every externally visible
operation must return a total code, with an explicit failure flag where
an input condition fails.  In particular, division by zero, a zero
projective tuple, and a malformed certificate must not produce an
unproved obligation in the implementation.  Define bounded division,
bounded gcd search, exact quotient tests, integer square-root search,
finite products, and coordinate scans as explicit combinations of
\texttt{PR.comp} and \texttt{PR.prec}.

The required interface is not merely a Lean function
\texttt{Nat -> Nat}; it consists of a code in \texttt{PR k}, its defining
equations as terms, and derivations of the equations used by subsequent
proofs.  Each fixed algorithm has such a code.  No universal total
primitive recursive evaluator for arbitrary primitive recursive programs
is assumed.  The uniform algorithms used here evaluate bounded polynomial
data and finite certificates, which is a different and much weaker task.

\paragraph{2. Elementary integer arithmetic inside the theory.}
Create \path{DG/Arithmetic.lean}.  Derive addition, multiplication,
order, and signed arithmetic laws from the defining equations, using
\texttt{PRA.qfInd}.  Prove the bounded division invariants
$a=qb+r$ and $r<b$ for $b>0$, Euclid's divisibility lemma, and the gcd and
coprime square-product lemmas used in descent.  Prime factorization can
be implemented by trial division through the input bound; it needs no
unbounded least-prime operator.  Whenever a proof is phrased with bounded
universal quantification, replace it by the equality-to-$1$ of its
primitive recursive Boolean product before invoking induction.

Create \path{DG/PolynomialCertificates.lean} next.  A reflected
polynomial normalizer should return both a coefficient array and an
object-theory equality proof.  Its soundness is proved by structural
recursion on the polynomial expression and the already derived ring
laws.  This reconstructs the role of the checked Lean identities without
treating an ambient \texttt{Int} identity as an object-theory axiom.
The relatively small identities in \texttt{DiophantineCore.lean} are
regression targets for this compiler.

\paragraph{3. Bounded scans and projective normalization.}
Create \path{DG/FiniteSearch.lean} and \path{DG/Projective.lean}.
For witness search, keep a flag and a candidate index in the recursive
state.  The soundness invariant says that either the flag is false or the
candidate is in the scanned interval and satisfies the Boolean test.
For completeness, fix a proposed successful index $j$ as a numerical
parameter and induct on the scan bound $i$ in the quantifier-free
statement
\[
 j\leq i\ \wedge\ p(j)=1\quad\Longrightarrow\quad
 \operatorname{anyLE}(p,i)=1.
\]
Here $p$ denotes a fixed primitive recursive symbol, not a quantified
function variable of the object language.  Universal generalization over
$j$ is performed after the induction.  Alternatively, maintain the
Boolean product of all failed tests as one additional state component.

Prove canonical projective normalization and uniqueness by gcd
divisibility and a fixed sign convention.  The projective point code
then has literal equality, so searches do not need a quotient type in
the object theory.  Instantiate the scan theorem with the explicit box
size $(2B+1)^{n+1}$ to obtain the constructive Northcott enumeration,
membership correctness, and finite cardinality bound.

\paragraph{4. Heights and polynomial-map certificates.}
Create \path{DG/Height.lean} and \path{DG/Resultant.lean}.
The upper height estimate follows by summing bounded monomial estimates.
For the lower estimate, formalize the two homogeneous B\'ezout
identities supplied by the Sylvester adjugate.  A certificate checker
first compares polynomial coefficients.  A separate soundness theorem
then shows that the common divisor of the two evaluated forms divides
the resultant constant, and combines the certificate's coefficient
bounds with the triangle inequality.

All determinant computations can use finite sums over coded
permutations, with inversion parity computed by bounded loops.
The adjugate identity is a finite polynomial identity, so its proof can
be produced by the same certificate mechanism.  This route does not
require an imported, unbounded algebraic-closure or Nullstellensatz
theorem.  The nonzero determinant test is an integer comparison and the
certificate coefficients are obtained by a finite explicit algorithm.
For a theorem uniform in the matrix size, also formalize cofactor expansion
and the sign-reversing permutation pairing as bounded-sum identities with
quantifier-free induction. Checking a numerical certificate verifies an
instance; it must not be substituted for this uniform arithmetic lemma.

\paragraph{5. Conic charts.}
Create \path{DG/Conics.lean}.  The checked second-intersection identity
settles only incidence.  The remaining formal lemmas must establish that
the produced tuple is nonzero under the nonsingularity and direction
hypotheses, handle the tangent parameter, prove the inverse chart
equations, and establish projective uniqueness and the stated height
bounds.  Work first with cleared integer coordinates and determinants.
Introduce normalized rational/projective codes only after these
polynomial equalities and nonvanishing guards have been proved.
In this way every division in the geometric prose has an explicit
arithmetical justification.

\paragraph{6. Descent and the elliptic example.}
Create \path{DG/FermatDescent.lean} and
\path{DG/CongruentNumber.lean}. Encode the positive integer solution
data as one natural number and define the descent step by the bounded
gcd, quotient, parity, and square-root operations.  Check separately
that it preserves the solution equations and decreases the displayed
rank.  These are quantifier-free implications about fixed primitive
recursive functions.

The generic descent argument is then compiled from the invariant
\[
 \operatorname{Good}(s_k)=1\quad\wedge\quad
 \operatorname{rank}(s_k)+k\leq\operatorname{rank}(s_0),
 \qquad s_{k+1}=\operatorname{Step}(s_k),
\]
at $k=\operatorname{rank}(s_0)+1$.  This is the exact invariant of the
checked \texttt{noBooleanDescent} theorem.  It uses no minimal
counterexample principle and no induction over an unbounded existential
formula.  Finally compile the rational triangle--elliptic-curve
transformations, including the nonzero-denominator and sign conditions,
and deduce the complete rational-point list for $y^2=x^3-x$.

\paragraph{7. Arithmetic dynamics.}
Create \path{DG/Dynamics.lean}.  Starting from the certified two-sided
height estimate, formalize strict height growth outside the explicit
threshold.  Above that threshold, iterate for a specified number of
steps and use quantifier-free arithmetic invariants to rule out a repeat.
Inside the height box, use the complete projective list and a bounded
orbit table.  Equality and membership are Boolean tests on normalized
codes, so the finite repeat test is primitive recursive.  This proves
both soundness and completeness of the algorithm listing all rational
preperiodic points.  It is important to prove the escape lemma before
discarding an orbit that leaves the box; finite search by itself would
not justify that exclusion.

\paragraph{8. Rational roots, local certificates, and denominator geometry.}
Create \path{DG/RationalRoots.lean} and
\path{DG/LocalObstructions.lean}.  The former needs only denominator
clearing, coprimality of powers, and the bounded divisor enumerator.
The latter proves reduction of polynomial evaluation modulo a supplied
modulus and preservation of unimodularity of a primitive integer tuple.
Its finite negative residue certificate is a Boolean conjunction, not a
local-to-global existence principle.

Create \path{DG/WeierstrassDenominators.lean}.  Prove reduced-fraction
uniqueness and the equation $e^2=q^3$ first.  The bounded factorization
module gives the exponent equality $2v=3u$ for each prime, hence the
square/cube denominator form.  Compile the independent square-test
enumerator and prove its list equivalence to the height predicate.
Finally \path{DG/HeightDescent.lean} can express
Theorem~\ref{thm:generation} as a schema for fixed group operations with
supplied object-theory derivations of the group laws, representative
coverage, and reducing-selector properties.  Its two induction invariants
are the truncated height rank and equation~\eqref{eq:descent-expansion}.
This conditional module must retain those hypotheses rather than silently
instantiating them for an arbitrary elliptic curve.

\paragraph{9. Object-theory theorem objects and native transfer.}
Create \path{DG/Statements.lean}, \path{DG/Proofs.lean}, and
\path{DG/NativeTransfer.lean}.  Write each theorem as a closed formula
over the preceding primitive recursive symbols.  Instantiate the
derived arithmetic and search lemmas to obtain
$[]\vdash[\mathtt{PRAAx}]\varphi_G$.  Run an audit recording every use
of \texttt{PRA.qfInd} and the associated proof of
\texttt{Fm.QF}.  Only then apply \texttt{pra\_iff\_native}.
The final integration test should print the axioms of the object-theory
proof terms and ensure that no claimed geometric conclusion has entered
as an extra axiom or as an assumption to a conditional transfer theorem.

\subsection{Completion criteria}

The ordinary Lean arithmetic core is complete and compiled.  The
projective normalization, full conic bijection, Fermat descent step,
resultant-height argument, and preperiodic-point enumeration have full
human proofs in this article but are not fully formalized in Lean here.
The deeply embedded \iPRA\ derivations of the geometric theorems are
also not supplied.  Their completion requires the modules just listed,
with the explicit encodings and induction invariants given here; it does
not require a new mathematical finiteness theorem.

To reproduce the completed part on a machine with the pinned Lean
toolchain, run
\begin{verbatim}
cd lean
lake build
\end{verbatim}
or run \texttt{lean DiophantineCore.lean} directly.  No network fetch of
Mathlib is needed for this core.  Integration with the ProveIt proof
syntax should use the pinned repository snapshot and its own build
configuration, rather than copying the standalone build's dependency
assumptions to the entire repository.
